\begin{abstract}
UAV swarms can support low-altitude IoT services, but their links are affected by mobility, weather, and limited onboard resources. Existing works usually study caching, routing, deployment, or task collaboration separately. This paper presents \pchorch, a predictive computation, communication, and caching orchestration framework for heterogeneous UAV swarms under environmental link uncertainty. The framework uses an ANN-based link-risk predictor, cache-value-aware scheduling, and mobility-aware dwell control. The lightweight practical controller, \pclr, uses predicted link margin and outage risk together with cache value, energy cost, delay cost, load, and handover cost. We also study \psr, an internal stability-regularized variant, and \etp, an exploratory event-triggered ablation. On blind paper-lock seeds across seven regimes, \pclr\ ranks first against external baselines. Under combined stress, it improves the aggregate objective by 6.47\% over Reactive-3C, reduces average delay by 0.61\%, p95 delay by 2.14\%, handovers by 24.52\%, and outage by 9.14\%. These improvements over Reactive-3C are Holm-corrected significant for delay, p95 delay, handovers, and outage. The tradeoff is that \pclr\ uses 0.72\% more energy and has 1.05\% lower useful cache hit ratio than Reactive-3C. \psr\ achieves the best aggregate non-oracle objective, while \etp\ does not provide reliable additional gain. The results show that predictive 3C orchestration is useful, but it does not dominate every metric.
\end{abstract}

\begin{IEEEkeywords}
Low-altitude economy, UAV swarms, edge intelligence, 3C orchestration, caching, link-risk prediction, mobility-aware scheduling, Internet of Things.
\end{IEEEkeywords}
